\documentclass[11pt]{article}
\usepackage{preambulo}
\addbibresource{bibliografia.bib}

\begin{document}

\portada{Prueba de compilación del formato}

\tableofcontents
\newpage

\section{Propósito del documento}

Este archivo no constituye una entrega. Su función es verificar que el preámbulo y el flujo de
compilación automática operan correctamente antes de redactar los documentos definitivos. Para
ello se incorporan todos los elementos tipográficos que serán utilizados en los informes, de
manera que cualquier componente ausente en la distribución de \TeX Live se manifieste en esta
etapa y no durante la preparación de las entregas.

El archivo se elimina una vez que los tres documentos se encuentren redactados.

\section{Elementos verificados}

\subsection{Tipografía y separación silábica}

Precisión, imputación, codificación, estandarización, reducción de la dimensionalidad.
La correcta aplicación de acentos y la separación silábica automática del idioma español
constituyen un requisito para la legibilidad del documento final.

\subsection{Notación numérica}

El área bajo la curva de la regresión logística alcanzó un valor de $0.879$, mientras que el
método de vecinos más cercanos obtuvo $0.835$. La comparación entre ambos mediante la prueba
de McNemar arrojó $p = 0.0004$.

\subsection{Tablas}

La Tabla \ref{tab:resultados} presenta los resultados obtenidos sobre el conjunto de prueba.

\begin{table}[H]
\centering
\caption{Resultados de los modelos sobre el conjunto de prueba}
\label{tab:resultados}
\small
\begin{tabular}{p{4.5cm} r r r r}
\toprule
\textbf{Modelo} & \textbf{AUC} & \textbf{Exactitud} & \textbf{Precisión} & \textbf{Recall} \\
\midrule
Regresión logística   & 0.879 & 0.827 & 0.736 & 0.519 \\
Árbol de decisión     & 0.847 & 0.807 & 0.701 & 0.447 \\
Vecinos más cercanos  & 0.835 & 0.797 & 0.721 & 0.356 \\
\midrule
Referencia trivial    & 0.500 & 0.741 & 0.000 & 0.000 \\
\bottomrule
\end{tabular}
\end{table}

\subsection{Figuras}

La Figura \ref{fig:clases} muestra la distribución de la variable objetivo.

\begin{figure}[H]
    \centering
    \includegraphics[width=0.75\textwidth]{clases.png}
    \caption{Distribución de la clase objetivo en el conjunto de entrenamiento.}
    \label{fig:clases}
\end{figure}

\subsection{Fragmentos de código}

\begin{lstlisting}[caption={Ajuste del escalador únicamente con el conjunto de entrenamiento},label={lst:escalado}]
escalador = StandardScaler()
X_entrenamiento_escalado = escalador.fit_transform(X_entrenamiento)
X_prueba_escalado = escalador.transform(X_prueba)
\end{lstlisting}

\subsection{Expresiones matemáticas}

Los límites empleados para el acotamiento de valores atípicos se definen a partir del rango
intercuartílico como

\[
L_{\text{inf}} = Q_1 - 1.5 \cdot \text{RIC},
\qquad
L_{\text{sup}} = Q_3 + 1.5 \cdot \text{RIC}
\]

\subsection{Algoritmos}

\begin{algorithm}[H]
\caption{Acotamiento de valores atípicos}
\label{alg:acotar}
\begin{algorithmic}[1]
\Require Conjunto de entrenamiento $E$, conjunto de prueba $P$
\Ensure Ambos conjuntos con valores acotados
\For{cada columna numérica $c$}
    \State Calcular $Q_1$ y $Q_3$ sobre $E$
    \State $L_{\text{inf}} \gets Q_1 - 1.5 \cdot (Q_3 - Q_1)$
    \State $L_{\text{sup}} \gets Q_3 + 1.5 \cdot (Q_3 - Q_1)$
    \State Acotar $c$ en $E$ y en $P$ con los mismos límites
\EndFor
\end{algorithmic}
\end{algorithm}

\subsection{Citas y bibliografía}

El conjunto de datos proviene del repositorio UCI \parencite{uci_support2} y corresponde al
estudio clínico descrito por \textcite{knaus1995support}. Las pruebas estadísticas empleadas
para la comparación de clasificadores siguen los lineamientos de \textcite{demsar2006}. La
implementación se realizó con la biblioteca \textit{scikit-learn} \parencite{scikit-learn}.

\subsection{Enlaces}

El repositorio del proyecto se encuentra disponible en
\url{https://github.com/JozMdz/Realidad_Virtual}.

\section{Criterio de verificación}

La compilación se considera correcta si el documento generado incluye portada, índice,
numeración de página en el margen inferior derecho, tablas numeradas en romanos, figura,
fragmento de código, expresiones matemáticas, algoritmo y referencias en formato APA.

\printbibliography[title={Referencias}]

\end{document}
